% ============================================================ Capítulo 4
\chapter{Ingeniería de requisitos}\label{cap:requisitos}

\section{Criterio de clasificación}
Cada requisito se clasifica en una de tres categorías, para no atribuir a los documentos algo que no dicen:
\textbf{Documentado} (cita del taller o de una clase con página), \textbf{Diseño} (decisión técnica
justificada en un ADR) y \textbf{Extensión} (funcionalidad no exigida que aporta valor y se declara como tal).
La matriz se mantiene en \ruta{docs/requisitos/matriz_trazabilidad.csv} y se transforma en Markdown y LaTeX con
\cmd{uv run python scripts/generar\_matriz.py}, de modo que existe una única fuente de verdad.

\section{Requisitos funcionales y no funcionales}
\begin{table}[H]\centering\small
\caption{Resumen cuantitativo de la matriz}\label{tab:resumen-matriz}
\begin{tabular}{@{}lcccc@{}}\toprule
Categoría & Cantidad & Implementado* & Parcial & Documentado/pendiente\\\midrule
Funcionales del taller (RF) & 20 & 15 & 4 & 1\\
No funcionales (RNF) & 6 & 4 & 1 & 1\\
Extensiones (EXT) & 6 & 5 & 0 & 1\\\bottomrule
\end{tabular}
\end{table}
(*) Incluye filas implementadas con una verificación complementaria pendiente, indicada en la matriz. Los estados ``parcial'' tienen una causa común y explícita: dependen de una acción que debe ejecutarse en el
equipo  (scraping real, evaluación con un LLM real, publicación en GitHub, instalación de Engram/GGA).

\section{Casos de uso}
\begin{figure}[H]\centering
\begin{tikzpicture}[node distance=0.45cm and 1.2cm]
  \node[actor] (cli) {Cliente};
  \node[actor, below=3.2cm of cli] (dev) {Equipo\\desarrollo};
  \node[actor, below=2.6cm of dev] (ag) {Agente de\\código};
  \begin{scope}[every node/.style={ellipse,draw=fnaazul,fill=fnaazul!6,font=\sffamily\scriptsize,minimum width=3.7cm,minimum height=0.62cm}]
    \node[right=2.2cm of cli, yshift=1.1cm] (u4) {CU-04 Consultar Q\&A};
    \node[below=of u4] (u5) {CU-05 Generar FAQ};
    \node[below=of u5] (u6) {CU-06 Generar resumen};
    \node[below=of u6] (u1) {CU-01 Configurar fuentes};
    \node[below=of u1] (u2) {CU-02 Ejecutar scraping};
    \node[below=of u2] (u3) {CU-03 Construir base};
    \node[below=of u3] (u7) {CU-07 Comparar prompts};
    \node[below=of u7] (u8) {CU-08 Evaluar};
    \node[below=of u8] (u10) {CU-10 Auditar};
    \node[below=of u10] (u9) {CU-09 Consultar vía MCP};
  \end{scope}
  \begin{scope}[on background layer]
    \node[grupo=fnagris, fit=(u4)(u9), label={[font=\sffamily\scriptsize\bfseries]above:Sistema de conocimiento Fanalca}] {};
  \end{scope}
  \foreach \u in {u4,u5,u6} \draw[thick,fnagris] (cli) -- (\u.west);
  \foreach \u in {u1,u2,u3,u7,u8,u10} \draw[thick,fnagris] (dev) -- (\u.west);
  \draw[thick,fnagris] (ag) -- (u9.west);
  \draw[flechad] (u2.east) to[bend left=40] node[etq,right]{\guillemotleft precede\guillemotright} (u3.east);
\end{tikzpicture}
\caption{Diagrama de casos de uso}\label{fig:casos-uso}
\end{figure}

\section{Historias de usuario y criterios de aceptación}
\begin{table}[H]\centering\small
\caption{Historias de usuario principales con criterio verificable}\label{tab:historias}
\begin{tabularx}{\linewidth}{@{}l X X@{}}\toprule
ID & Historia & Criterio de aceptación (verificable)\\\midrule
HU-01 & Como cliente, se desea saber dónde queda la sede principal. & La respuesta contiene ``Acopi'' y ``Yumbo'',
cita \texttt{contacto\#000} y queda marcada como fundamentada.\\
HU-02 & Como cliente, se desea conocer el horario de atención. & Al no estar publicado, se devuelve la frase de
rechazo fija (P20).\\
HU-03 & Como analista, se requiere un archivo único con todo el conocimiento. & Existe
\ruta{data/gold/base_conocimiento.txt} con cada fragmento precedido por \texttt{[id] (sección | url)}.\\
HU-04 & Como ingeniero de datos, se requiere rastrear sin infringir \texttt{robots.txt}. & Ninguna petición HTTP
hacia rutas en \texttt{Disallow} (prueba \texttt{test\_rastreo\_completo\_respeta\_robots\_y\_limites}).\\
HU-05 & Como responsable de seguridad, se requiere bloquear intentos de manipulación. & ``Ignora todas las
instrucciones\ldots'' produce \texttt{modo\_contexto = bloqueado} sin invocar el modelo.\\
HU-06 & Como docente, se requiere ver las tres funcionalidades en la interfaz. & Pestañas Q\&A, Resumen y FAQ
operativas (pruebas e2e con \texttt{AppTest}).\\\bottomrule
\end{tabularx}
\end{table}

\section{Matriz de trazabilidad}
{\scriptsize
\begin{longtable}{@{}>{\raggedright\arraybackslash}P{1.0cm}>{\raggedright\arraybackslash}P{2.2cm}>{\raggedright\arraybackslash}P{4.0cm}>{\raggedright\arraybackslash}P{2.9cm}>{\raggedright\arraybackslash}P{2.8cm}>{\raggedright\arraybackslash}P{2.7cm}@{}}
\caption{Matriz de trazabilidad (generada desde el CSV)}\label{tab:matriz}\\\toprule
ID & Origen & Descripción & Componente & Prueba de aceptación & Estado\\\midrule\endfirsthead
\toprule ID & Origen & Descripción & Componente & Prueba de aceptación & Estado\\\midrule\endhead
RF-01 & Taller 1, p.1, \S{}1 Investigación & Investigar la empresa e identificar fuentes públicas (sitio oficial, redes sociales, medios) & docs/\allowbreak{}requisitos/\allowbreak{}analisis\_\allowbreak{}documentos.md, config.urls\_\allowbreak{}semilla & Revisión documental + test\_\allowbreak{}dominios\_\allowbreak{}permitidos & Parcial: redes sociales y medios registrados como fuentes manuales (sus términos de uso prohíben el scraping) \\ \hline
RF-02 & Taller 1, p.1, \S{}1 Definición del alcance & Definir el tipo de preguntas de un cliente en el primer contacto (horarios, sedes, productos, historia, contacto, procesos básicos) & evaluacion/\allowbreak{}preguntas.yaml (alcance) & test\_\allowbreak{}banco\_\allowbreak{}tiene\_\allowbreak{}al\_\allowbreak{}menos\_\allowbreak{}20\_\allowbreak{}preguntas\_\allowbreak{}distintas & Implementado \\ \hline
RF-03 & Taller 1, p.1, \S{}2 Extracción de datos & Web scraping del sitio oficial con librerías Python (requests, BeautifulSoup, Selenium) & scraping/\allowbreak{}rastreador.py, cliente\_\allowbreak{}http.py, extraccion.py, dinamico.py & test\_\allowbreak{}rastreo\_\allowbreak{}completo\_\allowbreak{}respeta\_\allowbreak{}robots\_\allowbreak{}y\_\allowbreak{}limites & Implementado y probado con sitio simulado, ejecución contra el sitio real pendiente en el computador local \\ \hline
RF-04 & Taller 1, p.1, \S{}2 Preprocesamiento & Limpiar el texto extraído (eliminar HTML, etc.) & pipelines/\allowbreak{}silver.py & test\_\allowbreak{}extraccion\_\allowbreak{}elimina\_\allowbreak{}ruido, test\_\allowbreak{}pipeline\_\allowbreak{}bronze\_\allowbreak{}silver\_\allowbreak{}gold\_\allowbreak{}desde\_\allowbreak{}rastreo & Implementado \\ \hline
RF-05 & Taller 1, p.1, \S{}2 Segmentación & Dividir en fragmentos pequeños y semánticamente coherentes (chunks) & pipelines/\allowbreak{}fragmentacion.py & test\_\allowbreak{}fragmentacion\_\allowbreak{}respeta\_\allowbreak{}encabezados\_\allowbreak{}y\_\allowbreak{}tamano & Implementado \\ \hline
RF-06 & Taller 1, p.1, \S{}2 Nota & Entregar el repositorio con todos los códigos y la base de conocimiento proveniente del scraping & repositorio completo, data/\allowbreak{}gold & Inspección del ZIP + uv run fabrica build --fuente scraping & Parcial: se entrega base semilla, la base de scraping se genera con un comando en el computador local \\ \hline
RF-07 & Taller 1, p.1, \S{}3 Selección del modelo & Elegir modelo open source vía Ollama o privado vía API & llm/\allowbreak{}proveedores.py & test\_\allowbreak{}crear\_\allowbreak{}modelo\_\allowbreak{}ollama\_\allowbreak{}configurado, test\_\allowbreak{}crear\_\allowbreak{}modelo\_\allowbreak{}sin\_\allowbreak{}clave\_\allowbreak{}falla & Implementado \\ \hline
RF-08 & Taller 1, p.1, \S{}3 Framework & Usar un framework LLM (LangChain recomendado o LlamaIndex) para orquestar & llm/\allowbreak{}cadenas.py (LCEL) & test\_\allowbreak{}qa\_\allowbreak{}v3\_\allowbreak{}con\_\allowbreak{}modelo\_\allowbreak{}simulado\_\allowbreak{}valida\_\allowbreak{}citas & Implementado \\ \hline
RF-09 & Taller 1, p.1, \S{}3 Prompt Engineering & Prompt robusto que responda basándose únicamente en el contexto proporcionado & llm/\allowbreak{}prompts/\allowbreak{}qa.yaml v2-v3, llm/\allowbreak{}salvaguardas.py & test\_\allowbreak{}prompts\_\allowbreak{}v2\_\allowbreak{}v3\_\allowbreak{}incluyen\_\allowbreak{}reglas\_\allowbreak{}anti\_\allowbreak{}alucinacion, test\_\allowbreak{}qa\_\allowbreak{}rechazo\_\allowbreak{}normalizado & Implementado \\ \hline
RF-10 & Taller 1, p.1, \S{}3 Experimentación & Experimentar con zero-shot, formato del prompt e instrucciones para evitar alucinaciones & catálogo v1/\allowbreak{}v2/\allowbreak{}v3, laboratorio de prompts, services/\allowbreak{}evaluacion.py & test\_\allowbreak{}catalogo\_\allowbreak{}tiene\_\allowbreak{}tres\_\allowbreak{}versiones\_\allowbreak{}y\_\allowbreak{}variables & Parcial: protocolo y herramientas listos, resultados con LLM real pendientes de ejecución local \\ \hline
RF-11 & Taller 1, p.1, \S{}3 Memoria & Consolidar todo el texto limpio en el prompt de sistema (memoria del asistente) & knowledge/\allowbreak{}base.py (modo completo) & test\_\allowbreak{}armar\_\allowbreak{}contexto\_\allowbreak{}modos & Implementado \\ \hline
RF-12 & Taller 1, p.1-2, \S{}4 Interfaz (texto parcialmente ilegible) & Interfaz de prueba (Streamlit/\allowbreak{}Gradio) que reciba una pregunta y muestre la respuesta generada & frontend/\allowbreak{}app.py & test\_\allowbreak{}qa\_\allowbreak{}desde\_\allowbreak{}boton\_\allowbreak{}de\_\allowbreak{}ejemplo & Implementado \\ \hline
RF-13 & Taller 1, p.2, \S{}5 Pruebas & Formular al menos 20 preguntas distintas para evaluar precisión y coherencia & evaluacion/\allowbreak{}preguntas.yaml (24), services/\allowbreak{}evaluacion.py & test\_\allowbreak{}evaluacion\_\allowbreak{}extractiva\_\allowbreak{}produce\_\allowbreak{}metricas & Implementado (métricas con LLM real pendientes) \\ \hline
RF-14 & Taller 1, p.2, \S{}5 Documentación & Informe con descripción del problema, planteamiento, preparación de datos, modelado y resultados & docs/\allowbreak{}latex (cap. Informe técnico del Módulo 1) & Compilación del PDF & Implementado (sección de resultados con LLM real pendiente) \\ \hline
RF-15 & Taller 1, p.2, \S{}5 Presentación & Sustentación en vivo de 15 min con demostración en Streamlit/\allowbreak{}Gradio, sin PowerPoint & docs/\allowbreak{}latex (guion de demostración) & Ensayo con guion & Documentado \\ \hline
RF-16 & Taller 1, p.2, Rúbrica (25 \%) & Implementación funcional de Resumen, FAQ y Q\&A usando LangChain & llm/\allowbreak{}cadenas.py, services/\allowbreak{}asistente.py & test\_\allowbreak{}faq\_\allowbreak{}y\_\allowbreak{}resumen, test\_\allowbreak{}qa\_\allowbreak{}contrato\_\allowbreak{}y\_\allowbreak{}auditoria & Implementado \\ \hline
RF-17 & Taller 1, p.2, Rúbrica (25 \%) & Consolidación del conocimiento en un archivo de texto limpio & data/\allowbreak{}gold/\allowbreak{}base\_\allowbreak{}conocimiento.txt & test\_\allowbreak{}construir\_\allowbreak{}gold\_\allowbreak{}genera\_\allowbreak{}artefactos & Implementado \\ \hline
RF-18 & Taller 1, p.2, Rúbrica (25 \%) & Documentación del proceso de experimentación con los prompts & llm/\allowbreak{}prompts/\allowbreak{}*.yaml (técnica, hipótesis), docs/\allowbreak{}prompts/\allowbreak{}bitacora\_\allowbreak{}experimentos.md & Revisión de bitácora & Parcial: bitácora con protocolo y resultados de línea base, resultados con LLM pendientes \\ \hline
RF-19 & Taller 1, p.2, Rúbrica (15 \%) & Dashboard funcional, bien presentado, que muestre las tres funcionalidades & frontend/\allowbreak{}app.py & test\_\allowbreak{}resumen\_\allowbreak{}y\_\allowbreak{}faq, test\_\allowbreak{}panel\_\allowbreak{}muestra\_\allowbreak{}metricas & Implementado \\ \hline
RF-20 & Taller 1, p.2, Rúbrica (15 \%) & Código limpio y bien organizado en GitHub & .github/\allowbreak{}workflows/\allowbreak{}ci.yml, ruff, mypy & uv run ruff check . \&\& uv run mypy & Implementado (publicación en GitHub: acción pendiente) \\ \hline
RNF-01 & Clase 2, p.2, \S{}3 & Gestionar el proyecto con uv (pyproject.toml, uv.lock), sin pip ni requirements.txt & pyproject.toml, uv.lock & CI: uv sync --locked & Implementado \\ \hline
RNF-02 & Clase 2, p.2, \S{}4 & Arquitectura Medallion Bronze (HTML + JSONL con URL, fecha, estado, hash), Silver (texto limpio deduplicado), Gold (Markdown + metadatos) & pipelines/\allowbreak{}bronze.py, silver.py, gold.py & test\_\allowbreak{}markdown\_\allowbreak{}gold\_\allowbreak{}con\_\allowbreak{}frontmatter, test\_\allowbreak{}reanudacion\_\allowbreak{}no\_\allowbreak{}repite\_\allowbreak{}capturas & Implementado \\ \hline
RNF-03 & Clase 2, p.3, Con criterio & Respetar robots.txt, User-Agent, timeout, pausas, solo datos públicos, Ley 1581 (Habeas Data) & scraping/\allowbreak{}robots.py, cliente\_\allowbreak{}http.py, pipelines/\allowbreak{}texto.py & test\_\allowbreak{}robots\_\allowbreak{}respeta\_\allowbreak{}disallow, test\_\allowbreak{}anonimizar\_\allowbreak{}correos\_\allowbreak{}conserva\_\allowbreak{}corporativos & Implementado \\ \hline
RNF-04 & Clase 3, \S{}2-\S{}5 & Entorno agéntico: OpenCode, MCP, skills, Engram, OpenSpec, GGA & opencode.json, .mcp.json, skills/\allowbreak{}, openspec/\allowbreak{} & openspec validate --all, test\_\allowbreak{}mcp & Parcial: configurado y validado, Engram y GGA requieren instalación local \\ \hline
RNF-05 & Clase 3, \S{}3 & Escanear skills de terceros con SkillSpector antes de instalar & skills/\allowbreak{}fna-seguridad & skillspector scan & Documentado \\ \hline
RNF-06 & Prompt maestro & Compatibilidad con Windows 11 (PowerShell) & scripts/\allowbreak{}*.ps1 & Análisis sintáctico con PowerShell 7.4 & Implementado (sintaxis verificada, ejecución en Windows pendiente) \\ \hline
EXT-01 & Decisión de diseño & Selección de contexto BM25 cuando la base excede el presupuesto & knowledge/\allowbreak{}recuperacion.py & test\_\allowbreak{}bm25\_\allowbreak{}prioriza\_\allowbreak{}fragmento\_\allowbreak{}correcto & Implementado \\ \hline
EXT-02 & Prompt maestro /\allowbreak{} Clase 1 (Fase 3) & API REST FastAPI con OpenAPI & api/\allowbreak{} & tests/\allowbreak{}integration/\allowbreak{}test\_\allowbreak{}api.py & Implementado \\ \hline
EXT-03 & Prompt maestro /\allowbreak{} Clase 3 & Servidor MCP propio de solo lectura sobre la base & mcp\_\allowbreak{}server/\allowbreak{}server.py & test\_\allowbreak{}sesion\_\allowbreak{}stdio\_\allowbreak{}real & Implementado \\ \hline
EXT-04 & Decisión de diseño & Auditoría de interacciones, pipeline y evaluaciones en SQLite con Alembic & repositories/\allowbreak{}, alembic/\allowbreak{} & test\_\allowbreak{}migracion\_\allowbreak{}alembic\_\allowbreak{}crea\_\allowbreak{}esquema & Implementado \\ \hline
EXT-05 & Decisión de diseño & Línea base extractiva sin LLM para pruebas y CI & llm/\allowbreak{}extractivo.py & test\_\allowbreak{}resumen\_\allowbreak{}con\_\allowbreak{}extractivo & Implementado \\ \hline
EXT-06 & Prompt maestro & Contenedores Docker (API, interfaz, Ollama) & infra/\allowbreak{}docker/\allowbreak{}Dockerfile, docker-compose.yml & docker compose up --build & Pendiente de validación (sin demonio Docker en el entorno de construcción) \\ \hline
V2-01 & Prompt maestro v2, \S{}11 & Obtener los datos por web scraping de autopartesfanalca.com.co y fanalca.com (o terceros citables) con el código incluido y un comando único & scraping/\allowbreak{}, pipelines/\allowbreak{}orquestador.py, cli.py (pipeline) & test\_\allowbreak{}metricas\_\allowbreak{}y\_\allowbreak{}graficas\_\allowbreak{}desde\_\allowbreak{}scraping & Implementado — ejecución real pendiente en Windows 11 (red del entorno de construcción bloqueada) \\ \hline
V2-02 & Prompt maestro v2, \S{}15 & Explicar cada etapa del proceso de datos con gráficas en el proyecto y el repositorio & pipelines/\allowbreak{}graficas.py, docs/\allowbreak{}graficas/\allowbreak{}, frontend/\allowbreak{}app.py (Proceso de datos) & test\_\allowbreak{}graficas\_\allowbreak{}desde\_\allowbreak{}semilla\_\allowbreak{}declaran\_\allowbreak{}ausencia\_\allowbreak{}de\_\allowbreak{}bronze, test\_\allowbreak{}pestanas\_\allowbreak{}v2\_\allowbreak{}proceso\_\allowbreak{}y\_\allowbreak{}simulador & Implementado (gráficas actuales con la base de demostración) \\ \hline
V2-03 & Prompt maestro v2, \S{}4 y \S{}29 & Usar Claude Code con el modelo adecuado por especialidad (Opus: arquitectura, diseño mecánico, restricciones, revisión — Sonnet: implementación) & .claude/\allowbreak{}agents/\allowbreak{}*.md (campo model), opencode.json, CLAUDE.md & Revisión de frontmatter & Implementado \\ \hline
V2-04 & Prompt maestro v2, cambios & Eliminar ingeniería inversa, JupyterLab, notebooks y Anaconda del proyecto & pyproject.toml, docs/\allowbreak{}, ADR-014 & Búsqueda de referencias residuales & Implementado \\ \hline
V2-05 & Prompt maestro v2, \S{}30–\S{}31 & Modelo de componentes con jerarquía y articulaciones fija, rotacional, prismática, esférica y compuesta & simulador/\allowbreak{}modelos.py, catalogo\_\allowbreak{}freno.py & test\_\allowbreak{}ensamblaje\_\allowbreak{}integro, test\_\allowbreak{}articulacion\_\allowbreak{}exige\_\allowbreak{}eje\_\allowbreak{}unitario\_\allowbreak{}y\_\allowbreak{}limites, test\_\allowbreak{}jerarquia\_\allowbreak{}cubre\_\allowbreak{}todas\_\allowbreak{}las\_\allowbreak{}piezas & Implementado \\ \hline
V2-06 & Prompt maestro v2, \S{}32 & Vista de despiece y reensamblaje reversible con secuencia de desmontaje & simulador/\allowbreak{}cinematica.py, web/\allowbreak{}src/\allowbreak{}paginas/\allowbreak{}MaquinaFreno.tsx & test\_\allowbreak{}despiece\_\allowbreak{}reversible\_\allowbreak{}y\_\allowbreak{}monotono, test\_\allowbreak{}secuencias\_\allowbreak{}desmontaje\_\allowbreak{}y\_\allowbreak{}reensamblaje & Implementado \\ \hline
V2-07 & Prompt maestro v2, \S{}33 & Freno de disco como caso principal con cálculo simplificado separado de la animación & simulador/\allowbreak{}fisica\_\allowbreak{}freno.py, web/\allowbreak{}src/\allowbreak{}modelos/\allowbreak{}freno.ts & test\_\allowbreak{}fisica\_\allowbreak{}formulas\_\allowbreak{}de\_\allowbreak{}referencia, modelos.test.ts & Implementado \\ \hline
V2-08 & Prompt maestro v2, \S{}34 & Generación procedural por CLI de Python y exportación GLB con manifiesto & simulador/\allowbreak{}generacion.py, web/\allowbreak{}public/\allowbreak{}modelos/\allowbreak{} & test\_\allowbreak{}generacion\_\allowbreak{}reproducible, test\_\allowbreak{}manifiesto\_\allowbreak{}completo\_\allowbreak{}y\_\allowbreak{}mallas\_\allowbreak{}cerradas & Implementado \\ \hline
V2-09 & Prompt maestro v2, \S{}35 & Visor web con órbita, zoom, selección, árbol, ficha, modos, velocidad, visibilidad, transparencia, sección, cámara, PNG y lecturas & web/\allowbreak{}src/\allowbreak{} & e2e/\allowbreak{}visor.spec.ts (freno 3D) & Implementado \\ \hline
V2-10 & Prompt maestro v2, \S{}36 & Integración RAG que separa información documentada y de demostración, sin procedimientos de reparación o seguridad & services/\allowbreak{}simulador.py, mcp\_\allowbreak{}server (consultar\_\allowbreak{}pieza) & test\_\allowbreak{}preguntas\_\allowbreak{}sobre\_\allowbreak{}piezas, test\_\allowbreak{}consultar\_\allowbreak{}pieza\_\allowbreak{}separa\_\allowbreak{}modelo\_\allowbreak{}y\_\allowbreak{}documentado & Implementado \\ \hline
V2-11 & Prompt maestro v2, \S{}37 & API REST del simulador con validación Pydantic de valores del cliente & api/\allowbreak{}routes/\allowbreak{}simulador.py & test\_\allowbreak{}cinematica\_\allowbreak{}valida\_\allowbreak{}limites, test\_\allowbreak{}estados\_\allowbreak{}validados\_\allowbreak{}y\_\allowbreak{}persistidos & Implementado \\ \hline
V2-12 & Prompt maestro v2, \S{}38 & Persistencia de estados y eventos en PostgreSQL con Alembic & repositories/\allowbreak{}, alembic/\allowbreak{}versions/\allowbreak{}0002\_\allowbreak{}simulador.py, docker-compose.yml & test\_\allowbreak{}migraciones\_\allowbreak{}alembic\_\allowbreak{}en\_\allowbreak{}postgresql, test\_\allowbreak{}repositorio\_\allowbreak{}en\_\allowbreak{}postgresql & Implementado \\ \hline
V2-13 & Prompt maestro v2, \S{}39 & Biblioteca de máquinas (Anatomy) con menú de tarjetas alimentado por el web scraping, recorrido < 30 s y modo experimento & services/\allowbreak{}maquinas.py, web/\allowbreak{}src/\allowbreak{}paginas/\allowbreak{} & test\_\allowbreak{}biblioteca\_\allowbreak{}de\_\allowbreak{}maquinas\_\allowbreak{}usa\_\allowbreak{}la\_\allowbreak{}base, interfaz.test.tsx, e2e/\allowbreak{}recorrido.spec.ts & Implementado \\ \hline
V2-14 & Prompt maestro v2, \S{}40 & Pruebas de jerarquía, IDs, transformaciones, límites, reversibilidad, carga del modelo, API, frontend y rendimiento & tests/\allowbreak{}, web/\allowbreak{}src/\allowbreak{}\_\allowbreak{}\_\allowbreak{}tests\_\allowbreak{}\_\allowbreak{}/\allowbreak{}, web/\allowbreak{}e2e/\allowbreak{} & pytest, Vitest, Playwright & Implementado \\ \hline
V2-15 & Prompt maestro v2, \S{}28 & Streamlit como interfaz principal que enlaza e integra el visor & frontend/\allowbreak{}app.py (Simulador 3D) & test\_\allowbreak{}pestanas\_\allowbreak{}v2\_\allowbreak{}proceso\_\allowbreak{}y\_\allowbreak{}simulador & Implementado \\ \hline

\end{longtable}}
La columna ``historia'' y la columna ``caso de uso'' del CSV completan la trazabilidad en la versión Markdown
(\ruta{docs/requisitos/matriz_trazabilidad.md}).
